\chapter{The Owed and the Read}
% OUTLINE Ch4 "The Owed and the Read" -- the honest close. What Vol 5 READ
% (the single invariant) vs what stays OWED (fluid face -> later volume; the
% residue's four resemblances earned as real curvature/Morse/squeeze/Lawvere ->
% DEVICE self-founding work, not prose; 128/48 + 18=2x9 -> device-lane); the
% series' final fence-inventory (points to the Fence Ledger for the table);
% ends where the series began (electron = name of a model). Bold-math/fenced-
% referent; device-lane owed items in prose, no identifiers; step-not-move.
% Gate: Kodo (earn) + Beastmaster (whole-book).

An honest instrument keeps two columns, and a book about an honest
instrument closes by balancing them. In one column is what the series
read; in the other, what it owes. The value of the reading is not lessened
by the size of the debt --- it is guaranteed by it, because a book that
owed nothing would be a book that claimed everything, and this series has
been, from its first page, an argument against exactly that.

What the series read is one number, and it is now on the page. The single
invariant --- the model's own inverse coupling, bracketed between
\textbf{129.6} and \textbf{137.7} in the model's own exact counts, its
deep display near 137.011 shown and unsigned --- was gathered off a chain
earned across four volumes, read through a finite gauge that halts by
construction, and stated with its fence riding in the symbol so that it
could never be mistaken for the world's. That is the whole of what this
series claims to have measured: one model's coupling, read to the model's
own resolution, neighboring the laboratory's value by a margin the book
finds remarkable and refuses to explain. A single number, honestly read.
It is a small claim, and its smallness is its strength.

What the series owes is longer, and the book names each debt in its own
column so no reader has to reconstruct the list.

The first debt is a whole face. This series has, since its founding text,
described its subject through four descriptions, and only three of them
were built: the geometry, the intuition, and the grammar of the earlier
volumes, with the measurement that grounds them. The fourth --- the fluid
description, where viscosity and dissipation live --- was named in every
volume and constructed in none. It is owed to a later volume, the one that
will build the fourth face as the middle volumes built the other three,
and until that volume exists the fluid face is a promise on the books, not
a reading in them. The third chapter refused to borrow against it; this
one records it unpaid.

The second debt lives in the device rather than in the prose, and honesty
requires saying which debts are of that kind. Across the middle volumes,
the residue of the instrument's self-describing loop was shown to
\emph{resemble} four great mathematical structures --- a curvature, a
topological obstruction, a caught limit, a logical fixed point --- and
each resemblance was fenced, every time, as an illustration of shape and
never an identity. Whether any of those resemblances can be \emph{earned}
--- whether the residue is, in the cited sense and not merely in shape, a
curvature or a fixed point --- is real work, and it is not the work of a
sentence. It belongs to the instrument's own foundations: to the program,
underway and not finished, by which the device grounds its own arithmetic
from within rather than resting on machinery handed to it. That is a debt
paid in the workshop, not on the page, and this book does not pay it here
or pretend it is paid. The same is true of the two configuring numbers the
first volume marked as chosen rather than derived, and of the finer
question of how the model's counted eighteen arises with the structure it
has: earnings owed to the device's own deepening, flagged in the volumes
that raised them and left, honestly, to be closed where such things are
closed.

Set the two columns side by side and the shape of the whole series stands
clear. It developed, in explicit mathematics, four descriptions of one
model of the electron --- geometry, intuition, grammar, and the
measurement beneath them --- and it read the single number those
descriptions converge on. It did not measure the physical electron; it
does not bridge to the laboratory's value; it holds no theory of fluids,
no claim about a second particle, no road to quantum gravity. Every one of
those is either a fence it keeps or a debt it names, and the difference
between a fence and a debt is the difference between a boundary that never
moves and a boundary a later volume may. The lab-gap is a fence; the fluid
face is a debt. A reader who wants the full inventory --- every fence held
and everything held back behind it --- will find it in the ledger that
closes the book, where a series whose spine has been its boundaries
finishes by tabulating them.

And so the series ends where it began. Its first volume opened by saying
that the electron, in these books, is the name of a mathematical model,
and that the argument would touch descriptions and never the particle in
the world. Five volumes later, that is exactly what it has done: built the
descriptions, read the model's number, held the fence at the one point
where holding it was hardest, and reached no further. The innermost skin
of the onion is the same vow as the outermost, kept the whole way down.
The number was read; the line was not crossed; and what remains owed is
written in the open, which is the only place this series has ever been
willing to leave anything.
